\documentclass[12pt, a4paper]{article}
\usepackage[utf8]{inputenc}
\usepackage[T1]{fontenc}
\usepackage[a4paper, top=1in, bottom=1in, left=1.1in, right=1.1in]{geometry}
\usepackage{lmodern}
\usepackage{microtype}
\usepackage{setspace}
\setstretch{1.2}
\usepackage{parskip}
\usepackage{xcolor}
\usepackage{tcolorbox}

% Colors
\definecolor{axonpurple}{RGB}{108,99,255}
\definecolor{axondark}{RGB}{18,18,18}
\definecolor{axongold}{RGB}{198,138,0}

\newtcolorbox{questionbox}[1][]{
  enhanced, breakable, colback=axonpurple!5, colframe=axonpurple!80,
  title=#1, fonttitle=\bfseries\large, coltitle=white,
  left=10pt, right=10pt, top=10pt, bottom=10pt
}

\newtcolorbox{answerbox}[1][]{
  enhanced, breakable, colback=black!2, colframe=black!30,
  title=#1, fonttitle=\bfseries\large, coltitle=black,
  left=10pt, right=10pt, top=10pt, bottom=10pt
}

\title{\textbf{\color{axonpurple} Axon Flow: Project Q\&A (PQA)}}
\author{Architecture \& Design Discussions}
\date{\today}

\begin{document}

\maketitle

\section*{Architecture Q\&A}

\begin{questionbox}[Question 1: D3 vs React Flow]
So for making a deep nested hierarchy, is this step a must (like using D3 inside the pipeline)? Or is there any shortcut in React Flow?
\end{questionbox}

\begin{answerbox}[Answer]
The short answer is: \textbf{Yes, using an external layout engine like D3 (or Dagre/ELK) is completely necessary.} There is no shortcut built directly into React Flow for this.

Here is exactly why this step is required:

\subsubsection*{1. React Flow is Just a "Dumb" Canvas}
React Flow is incredibly powerful for rendering interactive nodes, handling drag-and-drop, and drawing smooth bezier curves between them. However, it is mathematically "dumb" when it comes to layout. 

If you pass a list of nodes to React Flow without specifying exact \texttt{x} and \texttt{y} coordinates, it will just pile every single node on top of each other at coordinate \texttt{(0, 0)}. React Flow has no built-in logic to say \textit{"this is a child, so put it 100 pixels to the right and 50 pixels down."}

\subsubsection*{2. Why \texttt{d3-hierarchy} is the Perfect Partner}
To make a mind map look like a mind map, you have to calculate the spacing perfectly so branches don't overlap. This requires complex math. 

Instead of writing that math from scratch, your project uses \texttt{d3-hierarchy}:
\begin{itemize}
    \item \textbf{\texttt{d3.stratify()}} does the job of figuring out who is the parent, who is the child, and who is the grandchild (creating the deep nesting in memory).
    \item \textbf{\texttt{d3.tree()}} uses advanced algorithms to calculate the perfect \texttt{x} and \texttt{y} coordinates for every node so they fan out beautifully without overlapping.
\end{itemize}

\subsubsection*{The Pipeline Summary}
Because of this separation of concerns, the industry-standard pipeline for a tool like yours is exactly what you have built:
\begin{enumerate}
    \item \textbf{MongoDB:} Stores flat, fast-to-query data.
    \item \textbf{D3 (The Brains):} Calculates the perfect mathematical layout (X/Y coordinates) in memory.
    \item \textbf{React Flow (The Muscle):} Takes those calculated X/Y coordinates and physically draws them on the screen, adding interaction.
\end{enumerate}
\end{answerbox}

\end{document}
